\documentclass[10pt,a4paper]{amsart}
\usepackage[margin=19mm]{geometry}
\usepackage{amsmath,amssymb,mathtools}
\usepackage[scr=rsfs,cal=euler]{mathalfa}
\usepackage{xfrac}
\usepackage[unicode,hidelinks,bookmarks=false]{hyperref}
\newcommand{\ie}{i.e.\ }
\newcommand{\cf}{cf.\ }
\newcommand{\resp}{respectively\ }
\newcommand{\Subsection}{\subsection}
\providecommand{\nobreakdash}{\nobreak\hbox{-}}
\setlength{\emergencystretch}{2em}
\multlinegap0pt
\usepackage{bm}
\usepackage{bbm}
\usepackage{dsfont}
\usepackage{xspace}
\usepackage{cancel}
\usepackage{mathtools}
\usepackage{amsthm}
\makeatletter
\@ifundefined{theorem}{\newtheorem{theorem}{Theorem}}{}
\@ifundefined{lemma}{\newtheorem{lemma}[theorem]{Lemma}}{}
\makeatother
\title[Unconditional convergence of general Fourier series]{Unconditional convergence of general Fourier series}
\author{Vakhtang Tsagareishvili, Giorgi Tutberidze, Giorgi Cagareishvili}
\date{Source: arXiv:2310.02194v2 (CC BY 4.0)}
\begin{document}
\maketitle

\noindent\emph{Source and licence.} Unconditional convergence of general Fourier series. Version arXiv:2310.02194v2 (CC BY 4.0), distributed under the Creative Commons Attribution 4.0 International licence, as recorded in the arXiv licence metadata for this exact version and shown on the abstract page https://arxiv.org/abs/2310.02194v2. Full rights notice: licenses/arxiv-2310.02194-CC-BY-4.0.txt.\par\smallskip

\noindent\textit{Editorial scope.} Counted unit: the Theorem whose source label is \texttt{theorem2} in the article (source0.tex lines 253--263, source label \texttt{theorem2}), delivered together with the article's own complete original proof of it (source0.tex lines 265--398, one \texttt{proof} environment of 134 lines). No mathematics is written, repaired or re-derived: statement and proof are the article's own text line for line. Required context retained uncounted: eq7 (lines 111--114); lemma2 (lines 157--167); eq6 (lines 209--214). Every definition, display and label of those lines is retained unchanged. Omitted from this package and not redistributed: 1--110, 115--156, 168--208, 215--252, 264--264, 399--1024. Nothing inside a retained excerpt is omitted or altered: every retained body is delivered exactly as the article's lines are written, so no omission receipt of any kind accompanies this package. Citations: the retained excerpts carry no citation command at all, so no bibliography entry of the article is transcribed and no reference is redistributed. Reference adaptation: none was needed and none was made; every reference inside the delivered unit resolves inside it, so no unresolved reference or citation is produced. Printed slips kept literal: no printed word, symbol, hypothesis or number of the counted statement or of its proof was altered. Layout and class: the article's own class is \texttt{publmathdeb} with packages including amsmath, amssymb; the wrapper is the lane's pinned \texttt{amsart} editorial wrapper at 10pt on A4, which restates the theorem-like environments the retained text needs and adds the article's own macro definitions that the retained text needs (none), with extra wrapper packages (bm, bbm, dsfont, xspace, cancel, mathtools). The delivered numbering is the unit's own. Assets: no retained span contains a figure environment, a graphics-inclusion command, a PDF-page inclusion, a table or a TikZ picture; no image file is copied into this package, the package declares no asset dependency and no image file is redistributed with it. Licence: version arXiv:2310.02194v2 of this article is distributed under the Creative Commons Attribution 4.0 International licence, as recorded in the arXiv licence metadata for this exact version and shown on the abstract page https://arxiv.org/abs/2310.02194v2. A full rights notice is delivered at licenses/arxiv-2310.02194-CC-BY-4.0.txt. Characters: every retained excerpt is pure ASCII, so the delivered engine, XeLaTeX with its default Latin Modern text font, reports no missing character.
	\begin{eqnarray}
		\label{eq7} B_n(f):=\sum_{k=1}^{n} C_k(f) a_k &= &\int_{0}^{1} f(x) \sum_{k=1}^{n} a_k \varphi_{k}(x) dx  \\
		&= & \int_{0}^{1} f(x) P_n (a, \varepsilon, x)dx. \notag 
	\end{eqnarray}

	\begin{lemma} \label{lemma2}
		Suppose that $(a_n)$ is any sequence of real numbers, and define 
		\[
		p(\varepsilon) = 2 - \varepsilon, \quad q(\varepsilon) = \frac{2 - \varepsilon}{1 - \varepsilon}, \quad \text{where } \varepsilon \in (0,1).
		\]
		Then,
		\[
		\frac{1}{n} \left( \sum_{k=1}^{n} a_k^2 \right)^{1/2} = O_\varepsilon(1).
		\]
		
	\end{lemma}

		\begin{eqnarray}
			\label{eq6}
			\int_{0}^{1} f(x) P_n(a, \varepsilon, x) dx &=& \sum_{i=1}^{n-1}\left(f\left(\frac{i}{n}\right)-f\left(\frac{i+1}{n}\right)\right)\int_{0}^{i/n}P_n(a, \varepsilon, x)dx \notag  \\
			&+& \sum_{i=1}^{n}\int_{\left(i-1\right)/n}^{i/n}\left(f\left(x\right) -f\left(\frac{i}{n}\right)\right)P_n(a, \varepsilon, x) dx     \notag \\
			&+& f\left(1\right)\int_{0}^{1}P_n(a, \varepsilon, x)dx=I_1 +I_2  +I_3.   
		\end{eqnarray}

	\begin{theorem}
		\label{theorem2}Let \( (\varphi_n) \) be an orthonormal system (ONS) on \([0,1]\), and suppose that \( \varepsilon \in (0,1) \) is an arbitrary number.  
		Assume that for some sequence \( (b_n) \in \ell_{q(\varepsilon)} \),  
		\begin{equation} \label{eq11}
			\limsup_{n \to +\infty} M_n(b, \varepsilon) = +\infty.
		\end{equation}
		Then there exists a function \( g \in \mathrm{Lip}_1 \) such that for any \( \varepsilon \in (0,1) \),  
		\[
		\sum_{n=1}^\infty |C_n(g)|^{2 - \varepsilon} = +\infty.
		\]	
	\end{theorem}

	\begin{proof}
		Firstly, we suppose that for any $\varepsilon \in (0,1)$
		\begin{eqnarray*}
			\limsup_{n \rightarrow +\infty} \left|\int_{0}^{1} P_n (b, \varepsilon, x) dx\right| =+ \infty.
		\end{eqnarray*} 
		Then if $l(x)=1,$ when $x\in [0,1]$ and $$C_n(l)=\int_{0}^{1}l(x) \varphi_{k}(x) dx \ \ \ \ \ (n=1,2, \dots),$$
		we will have	
		
		\begin{eqnarray*}
			&&\limsup_{n \rightarrow +\infty} \left|\sum_{k=1}^{n} b_k C_k (l)\right| = \limsup_{n \rightarrow +\infty} \left|\sum_{k=1}^{n} b_k \int_{0}^{1}l(x) \varphi_{k}(x) dx \right| \\
			&&=  \limsup_{n \rightarrow +\infty}  \left|\sum_{k=1}^{n} b_k \int_{0}^{1} \varphi_{k}(x) dx \right| = \limsup_{n \rightarrow +\infty} \left|\int_{0}^{1} P_n (b, \varepsilon, x) dx\right| =+ \infty.
		\end{eqnarray*}
		Thus, for some ($b_n) \in l_{q(\varepsilon)}$ the series	
		\begin{eqnarray*}
			\sum_{k=1}^{\infty} b_k C_k (l)
		\end{eqnarray*}
		is divergent. Consequently $(C_k (l)) \notin l_{p(\varepsilon)}$ or 
		\begin{eqnarray*}
			\sum_{n=1}^{\infty} \left|C_n(l)\right|^{2-\varepsilon} =+\infty.
		\end{eqnarray*}
		As $l \in Lip_1,$ then in such case Theorem \ref{theorem2} is valid.
		
		Now we suppose that for any $\varepsilon \in (0,1)$
		\begin{eqnarray}
			\label{eq12} \left|\int_{0}^{1} P_n (b, \varepsilon, x) dx\right| = O(1).
		\end{eqnarray}
		In \eqref{eq6} we substitute $P_n( a, \varepsilon, x)=P_n(b, \varepsilon,x)$ and $f(x)= g_n (x),$ where	
		\begin{eqnarray}
			\label{eq13} g_n(x) = \int_{0}^{x} sign \int_{0}^{t} P_n(b,  \varepsilon, u)du dt.
		\end{eqnarray}
		Now we have $(1<q(\varepsilon)<+\infty)$
		\begin{eqnarray}
			\label{eq14} \int_{0}^{1} g_n (x) P_n (b, \varepsilon, x) dx
			&=& \sum_{i=1}^{n-1}\left(g_n\left(\frac{i}{n}\right)-g_n\left(\frac{i+1}{n}\right)\right)\int_{0}^{i/n}P_n(b, \varepsilon, x)dx \notag   \\
			&+& \sum_{i=1}^{n}\int_{\left(i-1\right)/n}^{i/n}\left(g_n\left(x\right) -g_n\left(\frac{i}{n}\right)\right)P_n(b, \varepsilon, x) dx  \notag    \\
			&+& g_n\left(1\right)\int_{0}^{1}P_n(b, \varepsilon, x)dx=h_1 +h_2  +h_3.   
		\end{eqnarray}
		According to \eqref{eq13}, the Cauchy inequality and Lemma 2, we obtain the following:	
		\begin{eqnarray} \label{eq15}
			\left|h_2\right| &=& \left| \sum_{i=1}^{n}\int_{\left(i-1\right)/n}^{i/n}\left(g_n\left(x\right) -g_n\left(\frac{i}{n}\right)\right)P_n(b, \varepsilon, x) dx \right|   \\ 
			&& \leq \frac{1}{n} \int_{0}^{1} \left|P_n (b, \varepsilon, x)\right| \leq \frac{1}{n} \left(\int_{0}^{1} P_n^2 (b, \varepsilon, x)dx\right)^{1/2} \notag\\   
			&=& \frac{1}{n} \left(\sum_{k=1}^{n} b_k^2\right)^{1/2}  = O_{\varepsilon}(1). \notag
		\end{eqnarray}
		
		Next, by using \eqref{eq12} and \eqref{eq13} we receive 
		\begin{eqnarray}
			\label{eq16}
			\left|h_3\right| = O(1).
		\end{eqnarray}
		
		Now, by $E_n$ we denote the set of any numbers $i\in {1,2,\dots,n-1},$ for all of which, there exists a point $y\in \left[\frac{i}{n}, \frac{i+1}{n} \right),$   such that 
		
		\begin{eqnarray*}
			sign\int_{0}^{y} P_n(b, \varepsilon, x) dx \ne sign\int_{0}^{(i+1)/n} P_n(b, \varepsilon, x) dx.
		\end{eqnarray*}
		
		Then as the function $\int_{0}^{x} P_n(b, \varepsilon, t) dt$  is a continuous on  $\left[\frac{i}{n}, \frac{i+1}{n} \right),$  there exists a point $x_{in}\in \left[\frac{i}{n}, \frac{i+1}{n} \right),$ such that   
		$$\int_{0}^{x_{in}} P_n(b, \varepsilon, x) dx = 0.$$
		From here when $i\in E_n$ and as
		\begin{eqnarray*}
			\int_{0}^{i/n} P_n(b, \varepsilon, x) dx  = \int_{0}^{x_{in}} P_n(b, \varepsilon, x) dx-\int_{i/n}^{x_{in}} P_n(b, \varepsilon, x) dx = -\int_{i/n}^{x_{in}} P_n(b, \varepsilon, x) dx,
		\end{eqnarray*}
		Therefore, by Lemma \ref{lemma2}, we obtain  $(1<q(\varepsilon)<+\infty)$
		\begin{eqnarray}
			\label{eq17}
			\sum_{i\in E_n}\left|\int_{0}^{i/n} P_n(b, \varepsilon, x) dx\right| &=& \sum_{i\in E_n} \left|\int_{i/n}^{x_{in}} P_n(b, \varepsilon, x) dx\right|   \\
			&\leq &\int_{0}^{1} \left|P_n(b, \varepsilon, x) dx\right| \leq \left( \int_{0}^{1}P_n^2(b, \varepsilon, x) dx\right)^{1/2} \notag \\
			&=& \left(\sum_{k=1}^{n}b_k^2\right)^{1/2} =O_\varepsilon(1)\sqrt{n} . \notag
		\end{eqnarray}
		
		Now, suppose that \( F_n = [0,1] \setminus E_n \). Then, by equation~\eqref{eq13}, we have
		
		\begin{eqnarray*}
			g_n \left(\frac{i}{n}\right) -g_n \left(\frac{i+1}{n}\right) 
			&& = - \int_{i/n}^{(i+1)/n} sign \int_{0}^{x} P_n(b, \varepsilon,t) dt dx \int_{0}^{i/n}  P_n(b, \varepsilon,x) dx \\
			&&= - \frac{1}{n} \left|\int_{0}^{i/n}  P_n(b, \varepsilon,t)dx\right|
		\end{eqnarray*}
		and
		\begin{eqnarray}
			\label{eq18} 
			&&\left|\sum_{i\in F_n} \left(g_n \left(\frac{i}{n}\right) -g_n \left(\frac{i+1}{n}\right)\right)\int_{0}^{i/n}  P_n(b, \varepsilon,x)dx\right| \\
			&&= \frac{1}{n} \sum_{i\in F_n} \left| \int_{0}^{i/n} P_n(b, \varepsilon,x)dx\right|.  \notag 
		\end{eqnarray}
		Using equations~\eqref{eq17} and \eqref{eq18}, we obtain
		
		\begin{eqnarray}
			\label{eq19} \left|h_1\right| &=& \left| \sum_{i=1}^{n-1}\left(g_n\left(\frac{i}{n}\right)-g_n\left(\frac{i+1}{n}\right)\right)\int_{0}^{i/n}P_n(b, \varepsilon, x)dx  \right| \notag \\
			&=&  \left|\sum_{i\in F_n} \left(g_n \left(\frac{i}{n}\right) -g_n \left(\frac{i+1}{n}\right)\right)\int_{0}^{i/n}  P_n(b, \varepsilon,x)dx\right.  \\
			&+& \left.\sum_{i\in E_n} \left(g_n \left(\frac{i}{n}\right) -g_n \left(\frac{i+1}{n}\right)\right)\int_{0}^{i/n}  P_n(b, \varepsilon,x)dx\right| \notag \\
			&\geq& \left|\frac{1}{n} \sum_{i\in F_n} \left|\int_{0}^{i/n}  P_n(b, \varepsilon,x)dx\right| - \frac{1}{n} \sum_{i\in E_n} \left|\int_{0}^{i/n}  P_n(b, \varepsilon,x)dx\right| \right| \notag\\
			&=& \left|\frac{1}{n} \sum_{i=1}^{n-1} \left|\int_{0}^{i/n}  P_n(b, \varepsilon,x)dx\right| - \frac{2}{n} \sum_{i\in E_n} \left|\int_{0}^{i/n}  P_n(b, \varepsilon,x)dx\right| \right|  \notag \\
			&\geq& M_n(b,\varepsilon) - O\left(\frac{1}{\sqrt{n}}\right).\notag
		\end{eqnarray}
		
		Finally, from equality \eqref{eq14} and by account of \eqref{eq15}, \eqref{eq16} and \eqref{eq19} we obtain
		\begin{eqnarray*}
			\left|\int_{0}^{1} g_n(x) P_n(b, \varepsilon,x)dx\right| \geq M_n(b,\varepsilon) - O\left(\frac{1}{\sqrt{n}}\right)-O(1).
		\end{eqnarray*}
		Thus, by the condition of the Theorem \ref{theorem2} we get
		\begin{eqnarray} \label{eq20}
			&&\limsup_{n \rightarrow +\infty}\left|\int_{0}^{1} g_n(x) P_n(b, \varepsilon,x)dx\right| = \limsup_{n \rightarrow +\infty} M_n(b,\varepsilon) =+\infty. 
		\end{eqnarray}
		
		Now we must consider the sequence of bounded and continuous linear functionals on $Lip_1$ (see \eqref{eq7})
		\begin{eqnarray*}
			B_n(f):= \int_{0}^{1} f(x)  P_n(b, \varepsilon,x)dx.
		\end{eqnarray*}
		So, from \eqref{eq20} we obtain
		\begin{eqnarray*}
			\limsup_{n \rightarrow +\infty}\left|B_n \left(g_n\right)\right| =+ \infty.
		\end{eqnarray*}
		Since
		\begin{eqnarray*}
			\left|\left|g_n\right|\right|_{Lip_1} = 	\left|\left|g_n\right|\right|_C +\sup_{x,y \in \left[0,1\right]}  \frac{\left|g_n(x)-g_n(y)\right|}{\left|x-y\right|} \leq 2,
		\end{eqnarray*}
		by the Banach-Steinhaus theorem, there exists a function $g\in Lip_1$ such that
		\begin{eqnarray*}
			\limsup_{n \rightarrow +\infty}\left|B_n \left(g\right)\right| =+ \infty.
		\end{eqnarray*}
		
		Consequently
		\begin{eqnarray*}
			\limsup_{n \rightarrow +\infty}\left|\int_{0}^{1} g(x)  P_n(b, \varepsilon,x)dx\right| 
			&=&	\limsup_{n \rightarrow +\infty}\left|\sum_{k=1}^{n} b_k \int_{0}^{1} g(x)  \varphi_{k}(x)dx\right| \\ 
			&=& \limsup_{n \rightarrow +\infty}\left|\sum_{k=1}^{n} b_k C_k(g)\right| =+ \infty.
		\end{eqnarray*}
		Hence the series 
		$$\sum_{k=1}^{\infty} b_k C_k(g)$$
		is divergent for $(b_k)\in l_{q(\varepsilon)} .$ So, it is evident that  $(C_k (g))\notin l_{p(\varepsilon)}$ for any $\varepsilon \in (0,1)$ or
		\begin{eqnarray*}
			\sum_{k=1}^{\infty} \left|C_n(g)\right|^{2-\varepsilon} =+\infty.
		\end{eqnarray*}
		
	\end{proof}

\end{document}
